\def\ChapterCode{STW}
\chapter
[\CO{[STW]} Stoffwechselstörungen]
{Stoffwechselstörungen}
\ChapterIntro
{%
\Lang{Abkz.} metabolische Störungen.

Die relevanteste Stoffwechselstörung im Rettungsdienst ist der
\index{Diabetes mellitus}\T{Diabetes mellitus}. Dieser
stellt eine Störung des Kohlehydrat-(zucker-)stoffwechsels dar.
\EE{Zucker ist ein Kohlehydrat}. 
}
{%
\begin{description}
  \item[Maintainer:] Sebastian Gabriel
  \item[Autoren:] Diverse
  \item[Reviewer:] Standard-Reviewprozess
  \item[Version:] {\VarDocVersionTypeName} ({\VarDocVersionTypeDescription})
  \item[SHA1:] {\DocGitCommit}
\end{description}

{\TxTeilkapitelAASS}
}


\clearpage
\section{Blutzuckerstoffwechsel: \HeadingKeyword{Diabetes Mellitus}}
\label{topic:diabetes}
\label{topic:diabetes-mellitus}



\Text
{\TxBeschreibung: Blutzuckerstoffwechsel}
{\DefinitionInPara{Der \T{Diabetes mellitus} ist eine chronische
Stoffwechselerkrankung, bei der die Blutzuckerregulation gestört ist.} 
Durch die Erkrankung oder die unsachgemäße Behandlung kann es zu 
folgenden Notfällen kommen:

\begin{itemize}
  \item \EE{{Hyp\underline{\E{o}}glykämie}}: \enquote{Unterzucker}
  \item \EE{{Hyp\underline{\E{er}}glykämie}}: \enquote{Überzucker}
\end{itemize}}
{}
{}
{}
{}


\Table
{Interpretation des
BZ-Wertes: Richtwerte}%1 Titel
{Der Blutzuckerspiegel wird in \E{Milligramm pro Deziliter} 
(\MgDl\Z{, entspricht \MetricUnit{mg\,\PC}})\Z{, manchmal 
auch in \E{Millimol pro Liter} (\MetricUnit{\nicefrac{mmol}{\Liter}}),} 
gemessen. Der
Wert ist abhängig von der letzten Mahlzeit des Patienten.
Nach der Nahrungsaufnahme steigt der Blutzuckerspiegel stark
an und fällt dann wieder auf einen \enquote{nüchternen
Normalwert} ab. Der normale \EE{Nüchtern-Blutzuckerspiegel}
bewegt sich zwischen \EE{\Range{80}{100}\MgDl*} 
\Z{(\Range{4,4}{5,6}\,\MetricUnit{\nicefrac{mmol}{\Liter}})}. Im
Akutfall kann der Normalbereich jedoch großzügiger
angesetzt werden, da bei Werten zwischen 60 und 200\MgDl*
keine Notfallsymptome auftreten.
\cite{HeroldInnereMedizin2012}}%2 Beschreibung
{}%3 Label
{}%4 Quelle
{}%5 Lizenz
{%
\begin{tabular}{cccl}
\toprule
\multicolumn{3}{c}{\TabHeading{Wert [{\MgDl}]}}
&
\TabHeading{Beurteilung (nüchtern)}
\\\midrule
0 & -- & 20 & Lebensgefahr
\\\addlinespace
21 & -- & 60 & stark erniedrigt, Bewusstseinsstörungen
\\\addlinespace
61 & -- & 80 & erniedrigt, leichte Symptome
\\\addlinespace
81 & -- & 100 & Nüchtern-Blutzucker
\\\addlinespace
101 & -- & 200 & erhöht (nüchtern)
\\\addlinespace
201 & -- & 399 & stark erhöht
\\\addlinespace
400 & -- & ?? & sehr stark erhöht
\\\addlinespace\bottomrule
\end{tabular}
}%6 Tabelle
{}%7 Float

\Text{Insulin}
{%
\T{Insulin} ist für den Zuckerstoffwechsel wichtig, da es
dafür sorgt, dass Zucker (\T{Glukose}) aus dem Blut in die
Zellen aufgenommen wird. Es \EE{senkt daher den
Blutzuckerspiegel}. Insulin wird therapeutisch zur 
Behandlung des Diabetes mellitus angewendet.
Wenn man künstlich Insulin spritzt, besteht die Gefahr eines
Unterzuckers (zu viel Insulin gespritzt oder zu wenig
gegessen ${\rightarrow}$  Spritz-/ Essabstand)
er\-fra\-gen.
Hinsichtlich des Insulins sind zwei
Störungen häufig:
\begin{enumerate}
  \item Es wird \EE{zu wenig Insulin gebildet} (Bildungsort: Beta-Zellen des Pankreas).
  \item Der Körper hat einen \EE{gesteigerten Bedarf an Insulin},
  weil er schon zu sehr daran gewöhnt ist. Der Körper
  reagiert auf die \enquote{normale} Dosis nicht mehr so wie er
  sollte, er braucht größere Mengen Insulin um den gleichen
  Effekt wie ein Gesunder zu erreichen (\EE{Gewöhnungseffekt}),
  irgendwann kann aber der Körper nicht mehr Insulin
  produzieren.
\end{enumerate}
}{

}{}{}{}


\begin{Synopsis}
  \SynopsisItem{Insulin senkt den Blutzuckerspiegel.}
\end{Synopsis}




\TextSlide{Diabetes mellitus: Typen}
{%
Ausgehend von den das Insulin betreffenden Störungen 
unterscheidet man 2 Typen des Diabetes
mellitus:
\begin{itemize}
  \item \EE{Typ I}:  \Z{(IDDM,
  insuline dependent DM)} 
  
  Die Insulinproduktion ist gestört
  und der Körper produziert zu wenig Insulin. Der Patient
  muss als Ersatz künstliches Insulin spritzen.  Die Patienten
  sind eher jung (Jugendliche, junge Erwachsene) und
  sportlich.
  
  \Z{Ursache für die Störung ist eine Störung des
  Immunsystems, bei der es zu einer Zerstörung der
  Insulin-produzierenden ${\beta}$-Zellen des Pankreas
  kommt. Durch den Produktionsausfall kommt es zu einem
  absoluten Mangel.}
  
  \item \EE{Typ II}:
  \Z{(NIDDM, non insuline dependent DM)} 
  
  Durch eine
  dauerhaft zu hohe Zuckerzufuhr kommt es im Körper zu
  Gewöhnungseffekten und einem Wirkungsverlust des vom
  Körper produzierten Insulins \Z{(Insulinresistenz)}. Als
  Folge muss der Körper immer mehr Insulin produzieren um
  den gleichen Effekt zu erreichen. Wenn das
  Produktionslimit erreicht ist, kommt es zu einer Erhöhung
  des Blutzuckerspiegels und zum Typ-II-Diabetes.
  
  Der Patient mit Typ-II-Diabetes benötigt oft kein Insulin.
  Die Erkrankung kann zuerst mittels Lebensstiländerungen
  und \EE{Medikamenten} behandeln. Erst im Endstadium muss
  mitunter auch Insulin gespritzt werden.
  
  Der Typ-2-Diabetes betrifft eher Patienten im mittleren
  Alter und aufwärts, aufgrund der zunehmend fett- und
  kohlehydratreichen Ernährungsgewohnheiten werden die
  Patienten aber immer jünger.
\end{itemize}
}{
\begin{itemize}
  \item Typ I
  \begin{itemize}
    \item Insulinproduktion gestört
    \item Insulin muss gespritzt werden
    \item Eher Jugendliche und junge Erwachsene
  \end{itemize}
  \item Typ II
  \begin{itemize}
    \item Lebensstil, Ernährung
    \item Insulin-Gewöhnungseffekt
    \item Eher ab mittlerem Alter
    \item Lebensstiländerung, Medikamente, Insulin
  \end{itemize}
\end{itemize}
}{}{}{}


\TextSlide{Spätfolgen}
{%
Die eigentlichen Probleme der Zuckerkrankheit sind -- zumindest 
beim Typ-II-Diabetiker -- nicht ein akute Beschwerden, sondern 
die Langzeitfolgen der Erkrankung.
Ein hoher Blutzuckerspiegel über Jahre hinweg schädigt den gesamten
Körper. Es kommt zu:
\begin{itemize}
  \item \EE{Gefäßschäden}: 
  \begin{itemize}
    \item Gefäßverschlüsse (kleine und große Gefäße)
    \item \EE{Beine}: \enquote{Diabetischer Fuß}, infolge der
    Mangeldurchblutung in der Peripherie heilen auch kleine
    Wunden nicht mehr, in weiterer Folge muss immer weiter
    amputiert werden.
    \item \EE{Herz}: Die Herzkranzgefäße (Koronar\-arterien)
    sind ebenso betroffen. Diabetiker sind Risikopatienten
    für Herz-Kreislauf-Erkrankungen wie zB Herz\-infarkt.
    \item \EE{Hirn}: Wie beim Herzen kommt es zu
    Gefäßschäden, daher ist das Risiko für
    Schlaganfälle erhöht.
    \item \EE{Niere}: Die sehr kleinen Gefäße der Niere
    werden geschädigt. Im weiteren Verlauf stellt die Niere
    ihre Arbeit ein, es kommt zu einer
    \enquote{\T{chronischen Niereninsuffizienz}}, die im
    Endstadium mittels Blutwäsche (Dialyse) behandelt werden
    muss. Dialyse-Patienten machen einen großen Teil der
    Krankentransporte aus.
  \end{itemize}
  \item  \EE{Nervenschäden} \Z{(periphere Polyneuropathie)}:
  Gefühlslosigkeit vor allem in den Extremitäten, aber auch am
  gesamten Körper. Krankheiten, die sonst schmerzhaft sind,
  können beim Diabetiker aufgrund der Nervenschäden ohne
  Schmerzen ablaufen, z.\,B. ein schmerzloser, \EE{stummer Herzinfarkt}.
  \item  \EE{Augenschäden}: Bis hin zur Erblindung
\end{itemize}
}{
\begin{itemize}
  \item Gefäßschäden (Herz, Hirn, Nieren, Extremitäten, \ldots)
  \item Nervenschäden
  \item \Ca{Cave \enquote{Stummer Infarkt}!}
  \item Augenschäden
\end{itemize}
}{}{}{}


\begin{Danger}
  \item Achtung: \enquote{Stummer}, schmerzarmer/-loser Herzinfarkt beim Diabetiker!
\end{Danger}



% \clearpage
\subsection{Bei der \HeadingKeyword{Hypoglykämie} herrscht ein Blutzuckermangel}
\label{topic:hypoglykaemie}
\index{Hypoglykämie}



\TextSlide{{\TxBeschreibung} und Ursachen}
{%
\DefinitionInPara{Die \T{Hypoglykämie} ist eine maßgebliche 
Erniedrigung des Blutzuckerspiegels.}
Die häufigsten Ursachen für eine plötzliche Hypoglykämie
sind \EE{Diätfehler} und/oder eine \EE{Überdosierung von Insulin}
oder bestimmten anderen Diabetes-Medikamenten. 
Aufgrund
besonderer \EE{Anstrengungen} kann es zu einem gesteigerten
Zuckerbedarf des Körper kommen, wodurch der
Blutzuckerspiegel sinkt. Dies kommt häufig im Rahmen von
Begleiterkrankungen, wie z.\,B. grippalen Infekten, vor.
Auch die Alkohol-Intoxikation kann zu einer Hypoglykämie
führen\opt{NIVDOC}{ (Das beim Ethanolabbau entstehende
\Abk{NADH} hemmt die
Glukoneogenese.)} \citep{LoefflerBiochemie7}. 
\bigskip
}{%
\begin{itemize}
  \item Zu viel Insulin (spritzen ohne zu essen)
  \item Fasten (spritzen ohne zu Essen)
  \item Besondere Anstrengungen, Begleiterkrankungen,
  insb. Infekte
  \item Alkohol-Intoxikation
\end{itemize}
}{}{}{}






\TextSlide{\TxSymptome}
{%
Der Patient mit Hypoglykämie ist meistens \EE{unruhig},
\EE{desorientiert} und \EE{agitiert}, oft auch unkooperativ und
aggressiv. Oft fällt ein Zittern und eine schweißige Haut
auf. Auch Kopfschmerzen sind häufig. Eine Tachykardie und
Tachypnoe ist häufig feststellbar. Manchmal klagt der
Patient über Übelkeit und Heißhunger.
Je tiefer der Blutzuckerspiegel sinkt, desto deutlicher
werden die \EE{Bewusstseinsstörungen}. Von der anfänglichen
Agitiertheit und Desorientiertheit kann es rasch zu einer
Somnolenz, zum Sopor und zum Koma kommen!

Eine Hypoglykämie kann   \EE{fast alle neurologischen  Störungen
vortäuschen}! Häufig sind Kopfschmerzen und
Bewusstseinsstörungen, evtl. auch Krampfanfälle. Bei jeder
neurologischen Symptomatik muss man daher an eine
Hypoglykämie denken!
}{
\begin{itemize}
  \item Unruhe, Zittern, Schweißausbruch
  \item Aggressives Verhalten, unkooperativ
  \item Neurologische Symptomatik
  \begin{itemize}
    \item Bewusstseinsstörungen, Verwirrtheit, Kopfschmerzen, Krampfanfälle,
    Bewusstlosigkeit
  \end{itemize}
  \item Heißhunger, Übelkeit.
  \item Tachykardie, Tachypnoe
  \item Haut feucht/blass
\end{itemize}
}{}{}{}


\begin{Danger}
  \item Bei der Hypoglykämie können Symptome auftreten, die
  den Patient wie alkoholisiert wirken lassen! \E{Hüte dich
  vor vorschnellen
  Diagnosen!}
\end{Danger}

\begin{Synopsis}
  \SynopsisItem{${\rightarrow}$ BZ-Messung bei jeder neurologischen Störung.}
\end{Synopsis}



% M %%%%%%%%%%%%%%%%%%%%%%%%%%%%%%%%%%%%%%%%%%%%%%%%%%% M %
% Hypoglykämie
\ProcedurePrintDefault{ME16020C}

% \clearpage % TEMP temp clearpage
\subsection{Bei der \HeadingKeyword{Hyperglykämie} ist der Blutzuckerspiegel erhöht}
\label{topic:hyperglykaemie}
\index{Hyperglykämie}




\Text{Hyperglykämie: Zustand oder Notfall?}
{%
\DefinitionInPara{Eine \T{Hyperglykämie} ist eine 
maßgebliche Erhöhung des Blutzuckerspiegels.}
Dies ist nicht automatisch
ein Notfall. Die Bedrohlichkeit ist davon
abhängig, wie rasch sich diese Erhöhung entwickelt hat. So
hat ein Patient, bei dem seit Jahren ein unerkannter
Diabetes mellitus besteht auch mit sehr hohen
Blutzuckerwerten (akut) keine Probleme. Der Körper hat sich
über die Zeit daran \E{gewöhnt}.

Ein anderer Patient, bei dem der Zustand plötzlich
eingetreten ist (z.\,B. Erstmanifestation eines Diabetes
mellitus Typ I bei einem 19-jährigen, oder wenn das Insulin
nicht gespritzt wurde, \ldots), kann mit dem gleichen
Blutzuckerwert bereits im Koma liegen.
}
{}
{}{}{}

\begin{Synopsis}
  \SynopsisItem{Beurteile nicht nur den Messwert, sondern immer auch den Zustand des
  Patienten.}
\end{Synopsis}

\TextSlide{Auswirkungen}
{
Im Grunde verursacht der Zuckerspiegel im Akutfall zwei
Probleme: Einerseits wirkt der Zucker wie ein Elektrolyt auf
den \EE{Flüssigkeitshaushalt} ein und stört diesen, es kommt
zu einer ungünstigen Wasserumverteilung. Andererseits
mangelt es den Körperzellen an Zucker, da das Insulin fehlt,
um den Zucker in die Zellen hineinzubekommen -- es entsteht ein
\EE{Nährstoffmangel}.
}{
\begin{itemize}
  \item Störung des Wasser- und Elektrolythaushaltes
  \begin{itemize}
    \item Umverteilung
    \item Zucker wie Elektrolyt
  \end{itemize}
  \item Nährstoffmangel in Zellen
\end{itemize}
}{}{}{}

\SlideCompensation{1}

\TextSlide{Ursachen}
{%
Die Hyperglykämie betrifft praktisch nur Diabetes mellitus-Patienten. Bei
Patienten, die einen noch unerkannten DM~Typ~I haben, ist
ein hyperglykämischer Notfall oft die Erstmanifestation (die
Diagnose Diabetes ist noch nicht bekannt!). Sonst sind
Diätfehler oder eine mangelnde Insulinzufuhr häufige
Ursachen. Auch Stress und Begleiterkrankungen wie z.\,B.
akute Infektionen kommen in Betracht, da hier der
Insulinbedarf steigt.
}{
\begin{itemize}
  \item Erstmanifestation eines DM~Typ~I
  \item Diätfehler
  \item Insulinzufuhr $\downarrow$
  \item Stress
  \item Infektionen (Insulinbedarf ${\uparrow}$)
\end{itemize}
}{}{}{}









\TextSlide{\TxSymptome}
{%
Der Patient hat i.\,d.\,R. ein erhöhtes Durstgefühl und
gleichzeitig eine \EE{vermehrte Harnproduktion} (Zucker wird
ausgeschieden). Er klagt über Abgeschlagenheit, in weiterer
Folge kann es zu \EE{Bewusstseinsstörungen} bis hin zum
\EE{Koma} (s.u.) kommen. Die Geschwindigkeit der
Veschlechterung ist je nach Patient sehr unterschiedlich.

\subparagraph{Koma} 
\label{topic:koma-diabetisches}
\label{topic:diabetisches-koma}
\Z{Wie oben bereits beschrieben wurde, kann es sowohl zu
einem Problem des Flüssigkeitshaushaltes und/oder der
Nährstoffversorung der Zellen kommen. Je nachdem welches
Problem im Vordergrund steht, kommt es zu \E{unterschiedlichen Symptomen}.}


Ist vorwiegend der Flüssigkeitshaushalt gestört kommt es zum
\Z{sog. hyperosmolaren} Koma, dessen Leitsymptom
erwartungsgemäß die Bewusstseinsstörung ist.
Steht der Nährstoffmangel im Vordergrund passiert etwas
Interessantes: Der Körper produziert sog. \Z{Ketonkörper}. Da
diese \E{sauer} sind, kommt es zu einer stoffwechselbedingten
\EE{Übersäuerung} (metabolische
\EE{Azidose}\ZFootnote{Genauer:
\I[Ketoazidose!diabetische]{diabetischen Ketoazidose}}
\FullPrettyRef{topic:azidose-metabolische}). 
Er \EE{hyperventiliert}, um \ce{CO2} abzuatmen und somit
den pH wieder zu normalisieren (Puffersystem). Es ist daher
bei diesen Patienten eine \EE{schnelle, tiefe Atmung}
feststellbar (\Z{\I{Kußmaul-Atmung}}).

Mitunter kann auch ein blumig/fruchtiger oder Azeton-artiger
Atemgeruch wahrgenommen werden. Dieser darf nicht mit einer
\enquote{Alkoholfahne} verwechselt werden!
}{
\begin{itemize}
  \item Durstgefühl
  \item Vermehrte Harnproduktion
  \item Abgeschlagenheit, Bewusstseinstrübung bis
  \item Bewusstlosigkeit (Koma)
  \begin{itemize}
    \item Problem Flüssigkeitshaushalt
    \begin{itemize}
      \item \Z{Hyperosmolares Koma}
    \end{itemize}
    \item Problem Nährstoffmangel
    \begin{itemize}
      \item \Z{Ketoazidotisches Koma}
      \item Azidose → Hyperventilation zur Kompensation
    \end{itemize}
  \end{itemize}
  \item Azeton-Geruch (Nicht mit Alkoholfahne verwechseln!)
\end{itemize}
}{}{}{}



% M %%%%%%%%%%%%%%%%%%%%%%%%%%%%%%%%%%%%%%%%%%%%%%%%%%% M %
% Hyperglykämisches Koma
\ProcedurePrintDefault{ME14012C}





%%%%%%%%%%%%%%%%%%%%%%%%%%%%%%%%%%%%%%%%%%%%%%%%%%%%%%%%%%%
%%%%%%%%%%%%%%%%%%%%%%%%%%%%%%%%%%%%%%%%%%%%%%%%%%%%%%%%%%%
%%%%%%%%%%%%%%%%%%%%%%%%%%%%%%%%%%%%%%%%%%%%%%%%%%%%%%%%%%%


\ChapterEnd